\chapter{Conditional Expectation and Martingales}

\begin{center}\itshape
A monkey types random capital letters, one per second.\\
It needs on average $26^{11}+26^4+26$ seconds to type ABRACADABRA,\\
and a casino full of fair gamblers explains why.
\end{center}

In eighteenth-century France, a \emph{martingale} was a betting system: double your stake after every loss, so that the first win recovers all the
losses and gains one unit. It seems infallible, and for a player with unlimited credit and unlimited time it is: the first win comes with probability one.
Real players have limits, and the system ruined many of them. When mathematicians took the word over in the 1930s, they used it for the fortune of a
player in a fair game, and the first great theorem about martingales explains the catch: with bounded stakes or bounded time, no strategy can turn a
fair game into a favourable one (Exercise~\ref{ex:doubling} does the arithmetic).

A \emph{martingale} is the fortune of a player in a fair game: whatever has happened so far, the expected fortune after
the next round equals the present one. Sums of independent centred variables are martingales, but so are many
processes with strong dependence: the fraction of red balls in an urn that reinforces itself, a branching population
divided by its mean, the conditional expectation of any quantity as information accumulates. Martingale theory, created
by Joseph Doob around 1940 from ideas of Paul L\'evy and Jean Ville, is the most flexible tool of modern probability.

To define a fair game we need the expectation of a variable given the information available at a time, that is,
\emph{conditional expectation} given a $\sigma$-algebra. We build it as an orthogonal projection in the space $L^2$, after
proving that this space is complete, and then extend it to integrable variables. Then come the three great theorems of
Doob: \emph{optional stopping} (a fair game stays fair when you choose when to stop, within limits), the \emph{maximal
inequality} (a martingale cannot wander far without its final value noticing), and \emph{convergence} (a martingale that
stays bounded on average converges). Each comes with applications: the waiting time for a word, a second proof of the
gambler's ruin, P\'olya's urn, the growth of branching populations. Azuma's inequality closes the chapter: Hoeffding's
bound of Chapter~4 holds for martingales, which proves concentration for quantities that depend on many independent
inputs, none of them too much.

\begin{think}
In a casino with fair games, a player decides to stop as soon as she is one euro ahead. Since a fair random walk
returns to every level, she is sure to leave with a profit. Does this contradict the fairness of the games? What
would you need to forbid to restore it?
\end{think}

\section{Conditioning on a partition}

Let $A_1,A_2,\ldots$ be a finite or countable partition of $\Omega$ into events, and $\mathcal G$ the $\sigma$-algebra of all unions of blocks $A_k$. Knowing
$\mathcal G$ means knowing which block occurred. For an integrable $X$, the natural guess for $X$ given this information is its
average over the block that occurred:
\[
\mathbb E[X\mid\mathcal G]=\sum_{k:\,P(A_k)>0}\frac{\mathbb E[X\mathbf 1_{A_k}]}{P(A_k)}\,\mathbf 1_{A_k}.
\]
This is a random variable, constant on each block. For example, if $X$ is a die roll and $\mathcal G$ tells whether it is even, then
$\mathbb E[X\mid\mathcal G]$ is $4$ on $\{X\text{ even}\}$ and $3$ on $\{X\text{ odd}\}$. If $Y$ is a discrete random variable, the blocks $\{Y=y\}$ give $\mathbb E[X\mid Y]$, a
function of $Y$: the conditional expectations of Chapter~2, put together into one random variable.

Two properties characterise this variable $Y=\mathbb E[X\mid\mathcal G]$: it is constant on the blocks, and for every $A\in\mathcal G$,
\[
\mathbb E[Y\mathbf 1_A]=\mathbb E[X\mathbf 1_A].
\]
For $A=A_k$ both sides are $\mathbb E[X\mathbf 1_{A_k}]$, and a general $A$ is a countable union of blocks: sum over them, by dominated convergence
(the partial sums are bounded by $|X|$ on one side, and by $\sum_k\frac{|\mathbb E[X\mathbf 1_{A_k}]|}{P(A_k)}\mathbf 1_{A_k}$, of expectation at most $\mathbb E|X|$, on the other). On blocks
of probability $0$ the value of $Y$ is arbitrary; this harmless ambiguity will reappear as ``almost surely''.

There is a third description. If $\mathbb E[X^2]<\infty$, the constant $c$ that minimises $\mathbb E[(X-c)^2\mathbf 1_{A_k}]$ is the average $\frac{\mathbb E[X\mathbf 1_{A_k}]}{P(A_k)}$ (expand the square),
so $\mathbb E[X\mid\mathcal G]$ is the variable constant on the blocks that is closest to $X$ in mean square. For a general $\sigma$-algebra
$\mathcal G$, which need not come from a partition (the information ``the value of a Gaussian variable'' does not), the
first description makes no sense, but the other two do. We take the third as the definition, which requires some
geometry in a space of random variables.

\section{The space \texorpdfstring{$L^2$}{L2}}

Let $L^2$ be the set of random variables $X$ with $\mathbb E[X^2]<\infty$, two variables being identified when they are equal almost surely.
For $X,Y\in L^2$, $|XY|\le\frac12(X^2+Y^2)$ is integrable, and
\[
\langle X,Y\rangle=\mathbb E[XY],\qquad\|X\|=\sqrt{\mathbb E[X^2]}
\]
define an inner product and its norm: $\|X\|=0$ only if $X=0$ almost surely. Cauchy--Schwarz $|\mathbb E[XY]|\le\|X\|\,\|Y\|$ holds (Chapter~5), and
expanding $\|X+Y\|^2$ with it gives the triangle inequality $\|X+Y\|\le\|X\|+\|Y\|$. So $L^2$ is a Euclidean space as in Tome~0, Chapter~7, except
that its dimension is infinite. Convergence in $L^2$ (Chapter~6) is convergence for this norm. The decisive property of
$L^2$ is that it has no holes.

\begin{theorem}[Riesz--Fischer]\label{thm:complete}
$L^2$ is complete: if $(X_n)$ is a Cauchy sequence in $L^2$ (for every $\varepsilon>0$ there is $N$ with $\|X_m-X_n\|\le\varepsilon$ for $m,n\ge N$), then there is
$X\in L^2$ with $\|X_n-X\|\to0$. Moreover a subsequence of $(X_n)$ converges to $X$ almost surely.
\end{theorem}
\begin{proof}
Choose $n_1<n_2<\cdots$ such that $\|X_m-X_n\|\le2^{-k}$ for $m,n\ge n_k$. Let $G=\sum_k|X_{n_{k+1}}-X_{n_k}|$, a sum of non-negative terms. Its partial sums $G_K$ satisfy
$\|G_K\|\le\sum_{k<K}2^{-k}\le2$ by the triangle inequality, and $G_K^2$ increases to $G^2$, so $\mathbb E[G^2]\le4$ by monotone convergence. Hence $G<\infty$ almost
surely, and on this event the series $\sum_k(X_{n_{k+1}}-X_{n_k})$ converges absolutely: $X_{n_k}$ has a limit. Let $X$ be this limit where it exists, and $0$
elsewhere; it is a random variable (Proposition~\ref{prop:measurable}). For $n\ge n_j$, Fatou's lemma gives
\[
\mathbb E[(X_n-X)^2]\le\liminf_k\mathbb E[(X_n-X_{n_k})^2]\le4^{-j}.
\]
So $X_n-X\in L^2$, hence $X\in L^2$, and $\|X_n-X\|\to0$.
\end{proof}

A \emph{closed subspace} $H$ of $L^2$ is a linear subspace that contains the limits in $L^2$ of its convergent sequences.

\begin{theorem}[orthogonal projection]\label{thm:projection}
Let $H$ be a closed subspace of $L^2$ and $X\in L^2$. There is $Y\in H$, unique up to almost sure equality, such that
\[
\|X-Y\|=\min_{Z\in H}\|X-Z\|.
\]
It is characterised by: $Y\in H$ and $\mathbb E[(X-Y)Z]=0$ for every $Z\in H$.
\end{theorem}
\begin{proof}
Let $d$ be the infimum of $\|X-Z\|$ over $Z\in H$, and $Z_n\in H$ with $\|X-Z_n\|^2\to d^2$. The parallelogram identity $\|a+b\|^2+\|a-b\|^2=2\|a\|^2+2\|b\|^2$
with $a=X-Z_m$ and $b=X-Z_n$ gives, with $\delta_n=\|X-Z_n\|^2$ and since $V=\frac12(Z_m+Z_n)$ is in $H$,
\[
\begin{aligned}\|Z_n-Z_m\|^2&=2\delta_m+2\delta_n-4\|X-V\|^2\eqbr\le2\delta_m+2\delta_n-4d^2,\end{aligned}
\]
which tends to $0$. So $(Z_n)$ is a Cauchy sequence, which converges to some $Y$ (Theorem~\ref{thm:complete}), and $Y\in H$ since $H$ is
closed. By the triangle inequality $\|X-Y\|=\lim\|X-Z_n\|=d$.

If $Y$ is a minimiser and $Z\in H$, then for all real $t$, $\|X-Y-tZ\|^2\ge\|X-Y\|^2$, that is $t^2\|Z\|^2-2t\,\mathbb E[(X-Y)Z]\ge0$; this forces $\mathbb E[(X-Y)Z]=0$.
Conversely, if $Y'\in H$ and $X-Y'$ is orthogonal to $H$, then for $Z\in H$ Pythagoras gives $\|X-Z\|^2=\|X-Y'\|^2+\|Y'-Z\|^2\ge\|X-Y'\|^2$, so $Y'$ is a
minimiser. Finally, if $Y$ and $Y'$ both satisfy the orthogonality, then $\mathbb E[(Y'-Y)Z]=0$ for all $Z\in H$, and $Z=Y'-Y$ gives $Y=Y'$ almost surely.
\end{proof}

\begin{figure}[ht]
\centering
\fig{ch10-projection}
\caption{Conditional expectation as a projection: $\mathbb E[X\mid\mathcal G]$ is the point of the subspace $L^2(\mathcal G)$ closest to $X$, and
the error $X-\mathbb E[X\mid\mathcal G]$ is orthogonal to every $\mathcal G$-measurable variable.}
\label{fig:projection}
\end{figure}

\section{Conditional expectation}

Let $\mathcal G\subseteq\mathcal F$ be a $\sigma$-algebra, the information available. A random variable $Z$ is \emph{$\mathcal G$-measurable} if $\{Z\le t\}\in\mathcal G$ for every $t$:
its value is known once we know which events of $\mathcal G$ occurred. By Proposition~\ref{prop:measurable} applied on $(\Omega,\mathcal G)$, sums, products,
limits and continuous functions of $\mathcal G$-measurable variables are $\mathcal G$-measurable. Let $L^2(\mathcal G)$ be the set of $X\in L^2$ that are almost
surely equal to a $\mathcal G$-measurable variable.

\begin{lemma}\label{lem:closed}
$L^2(\mathcal G)$ is a closed subspace of $L^2$.
\end{lemma}
\begin{proof}
It is a linear subspace. Let $X_n\in L^2(\mathcal G)$ converge to $X$ in $L^2$, and let $X_n'$ be $\mathcal G$-measurable with $X_n'=X_n$ almost surely. Then
$X_n'\to X$ in $L^2$, and by Theorem~\ref{thm:complete} (whose limit is unique up to null sets) a subsequence $X_{n_k}'$ converges to $X$ almost surely.
The variable $X'=\limsup_kX'_{n_k}$ where it is finite, $0$ elsewhere, is $\mathcal G$-measurable and equal to $X$ almost surely.
\end{proof}

\begin{theorem}[conditional expectation]\label{thm:condexp}
Let $X$ be integrable. There is an integrable $\mathcal G$-measurable $Y$ such that
\[
\mathbb E[Y\mathbf 1_A]=\mathbb E[X\mathbf 1_A]\qquad\text{for every }A\in\mathcal G,
\]
and it is unique up to almost sure equality. It is written $\mathbb E[X\mid\mathcal G]$. If $X\in L^2$, it is the orthogonal projection of $X$ on
$L^2(\mathcal G)$. If $X\ge0$, then $\mathbb E[X\mid\mathcal G]\ge0$ almost surely.
\end{theorem}
\begin{proof}
\emph{Uniqueness.} If $Y$ and $Y'$ both qualify, $A=\{Y>Y'\}$ is in $\mathcal G$ and $\mathbb E[(Y-Y')\mathbf 1_A]=0$. The variable $(Y-Y')\mathbf 1_A$ is non-negative with
expectation $0$, so it is $0$ almost surely, that is $P(Y>Y')=0$. By symmetry $Y=Y'$ almost surely.

\emph{Square-integrable $X$.} Let $Y$ be a $\mathcal G$-measurable version of the projection of $X$ on $L^2(\mathcal G)$ (Lemma~\ref{lem:closed} and
Theorem~\ref{thm:projection}). For $A\in\mathcal G$, $\mathbf 1_A\in L^2(\mathcal G)$, so $\mathbb E[(X-Y)\mathbf 1_A]=0$. And $Y$ is integrable, as $|Y|\le\frac12(1+Y^2)$. If $X\ge0$, take
$A=\{Y<0\}$: then $\mathbb E[Y\mathbf 1_A]=\mathbb E[X\mathbf 1_A]\ge0$ while $Y\mathbf 1_A\le0$, so $Y\mathbf 1_A=0$ almost surely and $Y\ge0$ almost surely.

\emph{Integrable $X\ge0$.} Let $X_n=\min(X,n)\in L^2$ and $Y_n$ the variable just built for $X_n$. By uniqueness $Y_{n+1}-Y_n$ is the one for $X_{n+1}-X_n\ge0$,
so it is non-negative almost surely. Off a null set $N\in\mathcal G$, $(Y_n)$ is non-negative and non-decreasing; let $Y=\lim Y_n$ there and $Y=0$ on
$N$. For $A\in\mathcal G$, monotone convergence on both sides of $\mathbb E[Y_n\mathbf 1_A]=\mathbb E[X_n\mathbf 1_A]$ gives $\mathbb E[Y\mathbf 1_A]=\mathbb E[X\mathbf 1_A]$. With $A=\Omega$, $\mathbb E[Y]=\mathbb E[X]<\infty$,
so $Y$ is finite almost surely, and we may set $Y=0$ where it is infinite.

\emph{Integrable $X$.} Write $X=X^+-X^-$ and take the difference of the two variables just built.
\end{proof}

For a random variable or vector $V$, we write $\mathbb E[X\mid V]$ for $\mathbb E[X\mid\sigma(V)]$, where $\sigma(V)$ is the $\sigma$-algebra of the events $\{V\in B\}$
($B$ Borel). If $V$ is discrete, $\sigma(V)$ comes from the partition into the events $\{V=v\}$ and $\mathbb E[X\mid V]$ is the variable of Section~1, since that
variable has the defining property.

\begin{proposition}\label{prop:cerules}
Let $X$, $X'$ be integrable and $a$, $b$ real. Almost surely:
\begin{enumerate}[label=(\alph*)]
\item $\mathbb E[aX+bX'\mid\mathcal G]=a\,\mathbb E[X\mid\mathcal G]+b\,\mathbb E[X'\mid\mathcal G]$;
\item if $X\le X'$ then $\mathbb E[X\mid\mathcal G]\le\mathbb E[X'\mid\mathcal G]$; and $\bigl|\mathbb E[X\mid\mathcal G]\bigr|\le\mathbb E\bigl[|X|\,\big|\,\mathcal G\bigr]$;
\item $\mathbb E\bigl[\mathbb E[X\mid\mathcal G]\bigr]=\mathbb E[X]$;
\item if $X$ is $\mathcal G$-measurable, $\mathbb E[X\mid\mathcal G]=X$; if $X$ is independent of $\mathcal G$, $\mathbb E[X\mid\mathcal G]=\mathbb E[X]$;
\item (\emph{taking out what is known}) if $Z$ is $\mathcal G$-measurable and $ZX$ is integrable, $\mathbb E[ZX\mid\mathcal G]=Z\,\mathbb E[X\mid\mathcal G]$;
\item (\emph{tower property}) if $\mathcal H\subseteq\mathcal G$ is a $\sigma$-algebra, $\mathbb E\bigl[\mathbb E[X\mid\mathcal G]\,\big|\,\mathcal H\bigr]=\mathbb E[X\mid\mathcal H]$;
\item if $X\in L^2$, then $\mathbb E[X\mid\mathcal G]^2\le\mathbb E[X^2\mid\mathcal G]$, and in particular $\mathbb E[X\mid\mathcal G]\in L^2$.
\end{enumerate}
\end{proposition}
\begin{proof}
Each time we check that the proposed variable is $\mathcal G$-measurable and has the defining property, and conclude by uniqueness.
Write $Y=\mathbb E[X\mid\mathcal G]$.
(a) Linearity of $\mathbb E$. (b) $X'-X\ge0$, so $\mathbb E[X'-X\mid\mathcal G]\ge0$ by the theorem; apply this to $-|X|\le X\le|X|$ for the second statement.
(c) Take $A=\Omega$. (d) The first statement is the definition. For the second, ``independent of $\mathcal G$'' means that $X$ is independent of
$\mathbf 1_A$ for each $A\in\mathcal G$, so $\mathbb E[X\mathbf 1_A]=\mathbb E[X]P(A)$ by Proposition~\ref{prop:indepcdf}.

(e) If $Z=\mathbf 1_B$ with $B\in\mathcal G$, then $\mathbb E[\mathbf 1_BX\mathbf 1_A]=\mathbb E[X\mathbf 1_{A\cap B}]=\mathbb E[Y\mathbf 1_{A\cap B}]=\mathbb E[\mathbf 1_BY\mathbf 1_A]$; by linearity the formula holds
for simple $\mathcal G$-measurable $Z\ge0$. If $Z\ge0$ and $X\ge0$, then $Y\ge0$; take simple $Z_n$ increasing to $Z$ (as in Chapter~5) and apply
monotone convergence to both sides of $\mathbb E[Z_nX\mathbf 1_A]=\mathbb E[Z_nY\mathbf 1_A]$: so $\mathbb E[ZX\mathbf 1_A]=\mathbb E[ZY\mathbf 1_A]$, and with $A=\Omega$, $ZY$ is integrable.
In general split $Z=Z^+-Z^-$ and $X=X^+-X^-$: the four products $Z^\pm X^\pm$ are bounded by $|ZX|$, and the formula holds for each.

(f) For $A\in\mathcal H$, also $A\in\mathcal G$, so $\mathbb E[Y\mathbf 1_A]=\mathbb E[X\mathbf 1_A]=\mathbb E\bigl[\mathbb E[X\mid\mathcal H]\mathbf 1_A\bigr]$: the variable $\mathbb E[X\mid\mathcal H]$ has the defining
property of the conditional expectation of $Y$ given $\mathcal H$.

(g) For each rational $q$, $(X-q)^2\ge0$ gives $\mathbb E[(X-q)^2\mid\mathcal G]\ge0$, that is, by (a) and (d), $\mathbb E[X^2\mid\mathcal G]-2qY+q^2\ge0$ almost surely. Off a
null set this holds for all rational $q$ at once, hence by continuity for all real $q$, in particular for $q=Y(\omega)$: $\mathbb E[X^2\mid\mathcal G]\ge Y^2$.
Then $\mathbb E[Y^2]\le\mathbb E[X^2]$ by (c).
\end{proof}

Rule (e) says that what is known at the time of conditioning behaves as a constant; rule (f) says that the coarser
information wins. Here is a computation that uses only the definition and a symmetry.

\begin{example}\label{ex:exchange}
Let $X_1,\dots,X_n$ be independent with the same law, integrable, and $S=X_1+\dots+X_n$. Then $\mathbb E[X_1\mid S]=\frac Sn$. Indeed, for a Borel set $B$, the
vectors $(X_1,S)$ and $(X_i,S)$ have the same law (swapping $X_1$ and $X_i$ does not change the law of $(X_1,\dots,X_n)$, a product law), so
$\mathbb E[X_i\mathbf 1_{\{S\in B\}}]$ does not depend on $i$. Summing over $i$, $n\,\mathbb E[X_1\mathbf 1_{\{S\in B\}}]=\mathbb E[S\mathbf 1_{\{S\in B\}}]$, which is the defining property of $\frac Sn$.
Knowing only the total, each term is expected to be its fair share, whatever the law.
\end{example}

\section{Martingales}

\begin{definition}\label{def:martingale}
A \emph{filtration} is a non-decreasing sequence of $\sigma$-algebras $\mathcal F_0\subseteq\mathcal F_1\subseteq\cdots\subseteq\mathcal F$: $\mathcal F_n$ is the information available at time
$n$. A sequence $(X_n)_{n\ge0}$ is \emph{adapted} if each $X_n$ is $\mathcal F_n$-measurable. An adapted sequence of integrable variables is
a \emph{martingale} if
\[
\mathbb E[X_{n+1}\mid\mathcal F_n]=X_n\qquad\text{for all }n\ge0,
\]
a \emph{supermartingale} if $\mathbb E[X_{n+1}\mid\mathcal F_n]\le X_n$, and a \emph{submartingale} if $\mathbb E[X_{n+1}\mid\mathcal F_n]\ge X_n$.
\end{definition}

A supermartingale is an unfavourable game (the casino is right to like the word). By the tower property, a
martingale satisfies $\mathbb E[X_m\mid\mathcal F_n]=X_n$ for all $m\ge n$, and $\mathbb E[X_n]=\mathbb E[X_0]$ for all $n$. When no filtration is mentioned, $\mathcal F_n$ is
generated by the variables that drive the process up to time $n$.

\begin{example}\label{ex:martingales}
\begin{enumerate}[label=(\alph*)]
\item \emph{Sums.} If $\xi_1,\xi_2,\ldots$ are independent, integrable and centred, and $\mathcal F_n=\sigma(\xi_1,\dots,\xi_n)$, then $S_n=\xi_1+\dots+\xi_n$ is a
martingale: $\mathbb E[S_{n+1}\mid\mathcal F_n]=S_n+\mathbb E[\xi_{n+1}]$ by (a) and (d).
\item \emph{Squares.} If moreover $\mathbb E[\xi_k^2]=\sigma^2$ for all $k$, then $S_n^2-n\sigma^2$ is a martingale. Indeed
$S_{n+1}^2=S_n^2+2S_n\xi_{n+1}+\xi_{n+1}^2$, and $\mathbb E[S_n\xi_{n+1}\mid\mathcal F_n]=S_n\mathbb E[\xi_{n+1}]=0$ by (e) and (d), while $\mathbb E[\xi_{n+1}^2\mid\mathcal F_n]=\sigma^2$.
\item \emph{Products.} If $\xi_1,\xi_2,\ldots$ are independent, non-negative, with mean $1$, then $P_n=\xi_1\cdots\xi_n$ is a martingale. With $\xi_k$ equal to $0$
or $2$ with probability $\frac12$, $P_n$ is the fortune of a player who stakes everything on a fair coin at each round: $P_n\to0$ almost surely
while $\mathbb E[P_n]=1$.
\item \emph{Information.} If $X$ is integrable and $(\mathcal F_n)$ any filtration, $M_n=\mathbb E[X\mid\mathcal F_n]$ is a martingale, by the tower property: the
best guess of $X$ is revised fairly as information arrives.
\item \emph{Exponential martingale of a walk.} If $S_n$ is the walk with steps $+1$ with probability $p$ and $-1$ with probability $q=1-p$, then
$(\frac qp)^{S_n}$ is a martingale, as $p\frac qp+q\frac pq=1$; it is a product as in (c).
\end{enumerate}
\end{example}

A \emph{betting strategy} for a game with outcomes $X_k-X_{k-1}$ is a sequence $(H_k)_{k\ge1}$ where $H_k$, the stake on round $k$, is decided
before the round: $H_k$ is $\mathcal F_{k-1}$-measurable. Such a sequence is called \emph{predictable}. The player's gain after $n$ rounds is
\[
(H\cdot X)_n=\sum_{k=1}^nH_k(X_k-X_{k-1}),
\]
a discrete version of the stochastic integral of Chapter~11.

\begin{proposition}[no system beats a fair game]\label{prop:transform}
Let $(H_k)$ be predictable and bounded, $|H_k|\le K$. If $X$ is a martingale, $H\cdot X$ is a martingale. If $X$ is a supermartingale
and $0\le H_k\le K$, then $H\cdot X$ is a supermartingale; similarly for submartingales.
\end{proposition}
\begin{proof}
$(H\cdot X)_n$ is adapted, and integrable since it is bounded by $K\sum_{k\le n}|X_k-X_{k-1}|$. By rule (e),
$\mathbb E[(H\cdot X)_{n+1}-(H\cdot X)_n\mid\mathcal F_n]=H_{n+1}\,\mathbb E[X_{n+1}-X_n\mid\mathcal F_n]$, which is $0$ for a martingale, and $\le0$ for a supermartingale when $H_{n+1}\ge0$.
\end{proof}

\section{Stopping times and optional stopping}

\begin{definition}\label{def:stopping}
A \emph{stopping time} is a random variable $T$ with values in $\{0,1,2,\ldots\}\cup\{\infty\}$ such that $\{T=n\}\in\mathcal F_n$ for every $n$: whether to stop at
time $n$ is decided with the information available at time $n$.
\end{definition}

The first time an adapted process enters a Borel set $B$ is a stopping time, since $\{T\le n\}=\bigcup_{k\le n}\{X_k\in B\}\in\mathcal F_n$. The last time it is in $B$ is
in general not: deciding it requires the future. The process \emph{stopped} at $T$ is $X^T_n=X_{T\wedge n}$, where $T\wedge n=\min(T,n)$.

\begin{proposition}\label{prop:stopped}
If $X$ is a martingale (or a supermartingale) and $T$ a stopping time, the stopped process $X^T$ is a martingale (or a supermartingale).
In particular $\mathbb E[X_{T\wedge n}]=\mathbb E[X_0]$ for a martingale, and $\mathbb E[X_{T\wedge n}]\le\mathbb E[X_0]$ for a supermartingale.
\end{proposition}
\begin{proof}
$X_{T\wedge n}=X_0+\sum_{k=1}^n\mathbf 1_{\{T\ge k\}}(X_k-X_{k-1})$, and $H_k=\mathbf 1_{\{T\ge k\}}$ is predictable, since $\{T\ge k\}$ is the complement of $\{T\le k-1\}\in\mathcal F_{k-1}$.
Apply Proposition~\ref{prop:transform} with $0\le H_k\le1$.
\end{proof}

\begin{theorem}[optional stopping, Doob]\label{thm:ost}
Let $X$ be a martingale and $T$ a stopping time. In each of the following cases, $X_T$ is integrable and $\mathbb E[X_T]=\mathbb E[X_0]$:
\begin{enumerate}[label=(\alph*)]
\item $T\le N$ for a constant $N$;
\item $T<\infty$ almost surely, and $|X_{T\wedge n}|\le K$ for all $n$, for a constant $K$;
\item $\mathbb E[T]<\infty$, and $|X_n-X_{n-1}|\le c$ for all $n$, for a constant $c$.
\end{enumerate}
For a supermartingale, the same holds with $\mathbb E[X_T]\le\mathbb E[X_0]$.
\end{theorem}
\begin{proof}
By Proposition~\ref{prop:stopped}, $\mathbb E[X_{T\wedge n}]=\mathbb E[X_0]$ for every $n$ (or $\le$). (a) Take $n=N$. In (b) and (c), $T<\infty$ almost surely, so $X_{T\wedge n}\to X_T$
almost surely. In (b), dominated convergence with the bound $K$ concludes. In (c),
$|X_{T\wedge n}-X_0|\le\sum_{k\le T\wedge n}|X_k-X_{k-1}|\le cT$, and $cT+|X_0|$ is an integrable bound: dominated convergence again.
\end{proof}

The hypotheses cannot simply be dropped. For the simple random walk started at $0$, let $T$ be the first time it reaches $1$.
By the gambler's ruin (Theorem~\ref{thm:ruin}), the walk reaches $1$ before $-N$ with probability $\frac N{N+1}$, for every $N$; so $T<\infty$ almost surely.
Yet $S_T=1$, and $\mathbb E[S_T]=1\ne0$. The stopped walk is unbounded below, and by case (c) $\mathbb E[T]=\infty$. This answers the opening question:
the player who stops when one euro ahead needs unlimited credit and an infinite expected playing time; with either a
credit limit or a time limit, case (b) or (a) applies and the game is fair again.

Optional stopping gives a second proof of the symmetric gambler's ruin, without solving any recurrence.

\begin{proposition}\label{prop:ruinost}
Let $S$ be the simple random walk started at $a\in\{0,\dots,N\}$, and $T$ the first time it hits $0$ or $N$. Then $P(S_T=N)=\frac aN$ and $\mathbb E[T]=a(N-a)$.
\end{proposition}
\begin{proof}
By Lemma~\ref{lem:exit}, $T<\infty$ almost surely. The walk $S_n$ is a martingale (Example~\ref{ex:martingales}(a), shifted by $a$), and $|S_{T\wedge n}|\le N$; by case (b),
$a=\mathbb E[S_T]=N\,P(S_T=N)$. Next, $S_n^2-n$ is a martingale (Example~\ref{ex:martingales}(b): $(S_n-a)^2-n$ is one, and $S_n^2-n$ differs from it by $2a(S_n-a)+a^2$), so by Proposition~\ref{prop:stopped}
$\mathbb E[S_{T\wedge n}^2]-\mathbb E[T\wedge n]=a^2$. As $n\to\infty$, $\mathbb E[S^2_{T\wedge n}]\to\mathbb E[S_T^2]=N^2\cdot\frac aN$ by dominated convergence, and $\mathbb E[T\wedge n]\to\mathbb E[T]$ by monotone
convergence. So $\mathbb E[T]=aN-a^2$.
\end{proof}

Now the monkey. Fix an alphabet of $m$ letters and a word $w=w_1\cdots w_\ell$. Say that $k\in\{1,\dots,\ell\}$ is an \emph{overlap} of $w$ if the
first $k$ letters of $w$ are also its last $k$ letters: $w_1\cdots w_k=w_{\ell-k+1}\cdots w_\ell$. The word itself ($k=\ell$) is always an overlap.

\begin{theorem}[waiting for a word]\label{thm:abra}
Letters $U_1,U_2,\ldots$ are independent and uniform on an alphabet of $m$ letters. Let $T$ be the first $n$ such that $U_{n-\ell+1}\cdots U_n=w$. Then
\[
\mathbb E[T]=\sum_{k\text{ overlap of }w}m^k.
\]
\end{theorem}
\begin{proof}
\emph{The casino.} Just before each time $j=1,2,\ldots$ a new gambler arrives with $1$ euro. At time $j$ she stakes it on $U_j=w_1$, at fair odds:
the casino pays $m$ times the stake if she is right, and keeps it otherwise. If she wins, she stakes all her fortune on
$U_{j+1}=w_2$, and so on, until she loses everything or has won $\ell$ times; then she stops betting. Let $W^j_n$ be her fortune after time
$n$ ($W^j_n=1$ for $n<j$). Each $W^j$ is a martingale for $\mathcal F_n=\sigma(U_1,\dots,U_n)$: if she bets at time $n+1$ on the letter $w_r$ (an event and an index
known at time $n$), then $W^j_{n+1}=W^j_n\cdot m\mathbf 1_{\{U_{n+1}=w_r\}}$, and $\mathbb E[m\mathbf 1_{\{U_{n+1}=w_r\}}\mid\mathcal F_n]=1$ by rules (d) and (e); otherwise $W^j_{n+1}=W^j_n$.
So the total net gain of the gamblers,
\[
M_n=\sum_{j\le n}(W^j_n-1)=\Bigl(\sum_{j\le n}W^j_n\Bigr)-n,
\]
is a martingale with $M_0=0$. At most $\ell$ gamblers are betting at any time, each with a fortune at most $m^\ell$, so $|M_{n+1}-M_n|\le\ell m^\ell+1$.

\emph{Integrability of $T$.} Cut time into blocks of $\ell$ letters. Each block spells $w$ with probability $m^{-\ell}$, independently, so
$P(T>k\ell)\le(1-m^{-\ell})^k$ and $\mathbb E[T]=\sum_nP(T>n)\le\ell\sum_k(1-m^{-\ell})^k=\ell m^\ell$.

\emph{Stopping.} By Theorem~\ref{thm:ost}(c), $\mathbb E[M_T]=0$, that is $\mathbb E[T]=\mathbb E\bigl[\sum_{j\le T}W^j_T\bigr]$. At time $T$, the gambler who arrived at time $j=T-k+1$, for
$1\le k\le\ell$, has made $k$ bets on $w_1,\dots,w_k$ against the letters $U_j\cdots U_T=w_{\ell-k+1}\cdots w_\ell$; she has won them all, with fortune $m^k$, exactly
when $k$ is an overlap. A gambler who arrived at $j\le T-\ell$ is ruined: she cannot have won $\ell$ bets before time $T$ (the word would have
appeared earlier), so she lost one. Hence $\sum_{j\le T}W^j_T=\sum_{k\text{ overlap}}m^k$, a constant.
\end{proof}

ABRACADABRA has overlaps $1$ (A), $4$ (ABRA) and $11$, so the monkey needs on average $26^{11}+26^4+26=3\,670\,344\,487\,444\,778$ letters, about
$116$ million years at one letter per second. The word ABRACADABRX, with no overlap but itself, takes only $26^{11}$.
Yet both words appear with the same long-run frequency, $26^{-11}$ per position: occurrences of ABRACADABRA come in clumps
(ABRACADABRACADABRA contains two), so the gaps between clumps are longer.

\begin{beauty}
The proof never computes a probability. Instead of asking when the word appears, it asks who has won when it does, and
fairness converts the answer into an expectation. The same casino computes the expected waiting time for any pattern
in any sequence of independent symbols, a question of daily importance in DNA analysis and in string matching.
\end{beauty}

\section{Doob's maximal inequality}

Chebyshev's inequality bounds the probability that $S_n$ is large at the final time $n$. For a martingale, the same bound holds
for the maximum over all times up to $n$.

\begin{theorem}[Doob's maximal inequality]\label{thm:doob}
Let $X$ be a non-negative submartingale, $X^*_n=\max_{k\le n}X_k$ and $\lambda>0$. Then
\[
\lambda\,P(X_n^*\ge\lambda)\le\mathbb E\bigl[X_n\mathbf 1_{\{X^*_n\ge\lambda\}}\bigr]\le\mathbb E[X_n].
\]
If moreover each $X_n$ is in $L^2$, then $\mathbb E[(X^*_n)^2]\le4\,\mathbb E[X_n^2]$.
\end{theorem}
\begin{proof}
Let $T$ be the first $k$ with $X_k\ge\lambda$ (a stopping time, $T=\infty$ if there is none), and $A=\{X^*_n\ge\lambda\}=\{T\le n\}$. Then $X_n-X_{T\wedge n}=\sum_{k\le n}\mathbf 1_{\{T<k\}}(X_k-X_{k-1})$,
with $H_k=\mathbf 1_{\{T\le k-1\}}$ predictable and in $[0,1]$; by Proposition~\ref{prop:transform} this is a submartingale vanishing at $0$, so $\mathbb E[X_{T\wedge n}]\le\mathbb E[X_n]$.
On $A$, $X_{T\wedge n}=X_T\ge\lambda$; on $A^c$, $X_{T\wedge n}=X_n$. Hence
\[
\lambda P(A)+\mathbb E[X_n\mathbf 1_{A^c}]\le\mathbb E[X_{T\wedge n}]\le\mathbb E[X_n],
\]
and subtracting $\mathbb E[X_n\mathbf 1_{A^c}]$ gives the first inequality.

For the second, $Y=X^*_n$ is in $L^2$, since $Y^2\le\sum_{k\le n}X_k^2$. As $Y^2=\int_0^\infty2\lambda\mathbf 1_{\{Y\ge\lambda\}}\dd\lambda$, Tonelli's theorem (A2) and the first inequality give
\[
\begin{aligned}\mathbb E[Y^2]&=\int_0^\infty2\lambda P(Y\ge\lambda)\dd\lambda\eqbr\le\int_0^\infty2\,\mathbb E\bigl[X_n\mathbf 1_{\{Y\ge\lambda\}}\bigr]\dd\lambda\\&=2\,\mathbb E[X_nY]\le2\|X_n\|\,\|Y\|\end{aligned}
\]
by Tonelli again and Cauchy--Schwarz. If $\|Y\|>0$, divide by it: $\|Y\|\le2\|X_n\|$.
\end{proof}

\begin{corollary}[Kolmogorov's inequality]\label{cor:kolmogorov}
If $M$ is a martingale with $M_n\in L^2$, then for $\lambda>0$,
\[
P\Bigl(\max_{k\le n}|M_k|\ge\lambda\Bigr)\le\frac{\mathbb E[M_n^2]}{\lambda^2}.
\]
\end{corollary}
\begin{proof}
$M_k^2$ is a non-negative submartingale: $\mathbb E[M_{k+1}^2\mid\mathcal F_k]\ge\mathbb E[M_{k+1}\mid\mathcal F_k]^2=M_k^2$ by rule (g). Apply Theorem~\ref{thm:doob} to it, with $\lambda^2$.
\end{proof}

\section{Convergence in \texorpdfstring{$L^2$}{L2}}

\begin{lemma}\label{lem:orth}
Let $M$ be a martingale with $M_n\in L^2$ for all $n$, and $D_k=M_k-M_{k-1}$. Then $\mathbb E[D_kZ]=0$ for every $Z\in L^2$ that is $\mathcal F_{k-1}$-measurable. In particular the
increments are orthogonal to each other and to $M_0$, and
\[
\mathbb E[M_n^2]=\mathbb E[M_0^2]+\sum_{k=1}^n\mathbb E[D_k^2].
\]
\end{lemma}
\begin{proof}
$D_kZ$ is integrable by Cauchy--Schwarz, and by rules (c) and (e), $\mathbb E[D_kZ]=\mathbb E\bigl[Z\,\mathbb E[D_k\mid\mathcal F_{k-1}]\bigr]=0$. For $j<k$, $D_j$ and $M_0$ are $\mathcal F_{k-1}$-measurable.
Expand the square of $M_n=M_0+D_1+\dots+D_n$.
\end{proof}

\begin{theorem}[convergence in $L^2$]\label{thm:L2conv}
Let $M$ be a martingale bounded in $L^2$: $\sup_n\mathbb E[M_n^2]<\infty$. Then there is $M_\infty\in L^2$ such that $M_n\to M_\infty$ almost surely and in $L^2$, and
$M_n=\mathbb E[M_\infty\mid\mathcal F_n]$ for every $n$.
\end{theorem}
\begin{proof}
By Lemma~\ref{lem:orth}, $\sum_k\mathbb E[D_k^2]<\infty$; let $r_m=\sum_{k>m}\mathbb E[D_k^2]$, which tends to $0$. For $m<n$, $\|M_n-M_m\|^2=\sum_{m<k\le n}\mathbb E[D_k^2]\le r_m$, so $(M_n)$ is a
Cauchy sequence in $L^2$, with a limit $M_\infty$ by Theorem~\ref{thm:complete}.

\emph{Almost sure convergence.} For fixed $m$, $(M_{m+k}-M_m)_{k\ge0}$ is a martingale for the filtration $(\mathcal F_{m+k})_k$, so by Corollary~\ref{cor:kolmogorov},
$P(\max_{m\le k\le n}|M_k-M_m|\ge\varepsilon)\le\frac{r_m}{\varepsilon^2}$. These events increase with $n$; their union gives
\[
P\Bigl(\sup_{k\ge m}|M_k-M_m|\ge\varepsilon\Bigr)\le\frac{r_m}{\varepsilon^2}.
\]
Let $V_m=\sup_{j,k\ge m}|M_j-M_k|\le2\sup_{k\ge m}|M_k-M_m|$. It decreases with $m$ to a limit $V$, and $P(V\ge2\varepsilon)\le P(V_m\ge2\varepsilon)\le\frac{r_m}{\varepsilon^2}\to0$.
So $V=0$ almost surely (take $\varepsilon=\frac1j$ for all $j$): almost surely $(M_n)$ is a Cauchy sequence of reals, hence converges. Its limit is
$M_\infty$ almost surely, since a subsequence converges almost surely to $M_\infty$ (Theorem~\ref{thm:complete}).

\emph{Closure.} For $A\in\mathcal F_n$ and $m\ge n$, $\mathbb E[M_m\mathbf 1_A]=\mathbb E[M_n\mathbf 1_A]$, and $|\mathbb E[M_m\mathbf 1_A]-\mathbb E[M_\infty\mathbf 1_A]|\le\|M_m-M_\infty\|\to0$ by Cauchy--Schwarz.
So $\mathbb E[M_\infty\mathbf 1_A]=\mathbb E[M_n\mathbf 1_A]$.
\end{proof}

For example, if $\varepsilon_1,\varepsilon_2,\ldots$ are independent fair signs $\pm1$ and $\sum_na_n^2<\infty$, then $\sum_n\varepsilon_na_n$ converges almost surely: its partial
sums form a martingale with $\mathbb E[M_n^2]=\sum_{k\le n}a_k^2$. The harmonic series with random signs $\sum_n\frac{\varepsilon_n}n$ converges almost surely, although
$\sum_n\frac1n$ diverges.

\section{Upcrossings and almost sure convergence}

Boundedness in $L^2$ is a strong hypothesis. Doob found a way to prove almost sure convergence from boundedness in
$L^1$ alone, by counting oscillations. Fix $a<b$. An \emph{upcrossing} of $[a,b]$ by $X_0,\dots,X_n$ is a passage from a value $\le a$ to a later
value $\ge b$; let $U_n[a,b]$ be the number of completed upcrossings, the largest $k$ such that there are times
$s_1<t_1<s_2<t_2<\dots<s_k<t_k\le n$ with $X_{s_i}\le a$ and $X_{t_i}\ge b$.

\begin{lemma}[Doob's upcrossing inequality]\label{lem:upcrossing}
If $X$ is a supermartingale, then $(b-a)\,\mathbb E\bigl[U_n[a,b]\bigr]\le\mathbb E\bigl[(X_n-a)^-\bigr]$.
\end{lemma}
\begin{proof}
Consider the strategy ``buy low, sell high'': start holding one unit as soon as $X\le a$, hold it until $X\ge b$, then wait until
$X\le a$ again, and so on. Formally, $H_1=\mathbf 1_{\{X_0\le a\}}$, and for $k\ge1$, $H_{k+1}=1$ if either $H_k=1$ and $X_k<b$, or $H_k=0$ and $X_k\le a$; otherwise $H_{k+1}=0$.
By induction $H_k$ is $\mathcal F_{k-1}$-measurable, with values in $\{0,1\}$. Each completed upcrossing contributes at least $b-a$ to $(H\cdot X)_n$; a
holding period still open at time $n$, started at a time $s$ with $X_s\le a$, contributes $X_n-X_s\ge X_n-a\ge-(X_n-a)^-$. So
\[
(H\cdot X)_n\ge(b-a)\,U_n[a,b]-(X_n-a)^-.
\]
By Proposition~\ref{prop:transform}, $H\cdot X$ is a supermartingale vanishing at $0$, so $\mathbb E[(H\cdot X)_n]\le0$, and the inequality follows.
\end{proof}

\begin{figure}[ht]
\centering
\fig{ch10-upcross}
\caption{The upcrossings of $[a,b]$ by a path: the shaded periods are those where the strategy ``buy at or below $a$, sell at or
above $b$'' holds the asset. Each completed upcrossing earns at least $b-a$.}
\label{fig:upcross}
\end{figure}

\begin{theorem}[martingale convergence, Doob]\label{thm:supconv}
Let $X$ be a supermartingale with $\sup_n\mathbb E|X_n|<\infty$. Then $X_n$ converges almost surely to an integrable limit $X_\infty$. In particular,
every non-negative supermartingale, and every non-negative martingale, converges almost surely.
\end{theorem}
\begin{proof}
Let $K=\sup_n\mathbb E|X_n|$. As $n\to\infty$, $U_n[a,b]$ increases to the total number $U_\infty[a,b]$ of upcrossings, and $(X_n-a)^-\le|X_n|+|a|$, so by monotone
convergence and the lemma $(b-a)\,\mathbb E[U_\infty[a,b]]\le K+|a|$. Hence $U_\infty[a,b]<\infty$ almost surely. On the event $\{\liminf X_n<a<b<\limsup X_n\}$ the path
crosses $[a,b]$ upwards infinitely often, so this event is null. The union of these events over rational $a<b$ is null, and off it
$\liminf X_n=\limsup X_n$ in $[-\infty,+\infty]$. By Fatou's lemma $\mathbb E[\liminf|X_n|]\le K$, so the limit is finite almost surely and integrable. For a
non-negative supermartingale, $\mathbb E|X_n|=\mathbb E[X_n]\le\mathbb E[X_0]$.
\end{proof}

The limit need not be the limit in $L^1$: the fortune $P_n$ of the all-in player of Example~\ref{ex:martingales}(c) converges to $0$ while $\mathbb E[P_n]=1$.
The exact condition for convergence in $L^1$, \emph{uniform integrability}, is studied with the Lebesgue integral in Tome~6, and
uniformly integrable martingales in Tome~10.

\begin{theorem}[P\'olya's urn]\label{thm:polyaurn}
An urn contains one red and one blue ball. At each step a ball is drawn uniformly, and put back together with a new
ball of the same colour. Let $R_n$ be the number of red balls after $n$ steps and $M_n=\frac{R_n}{n+2}$ their proportion. Then $M_n$ is a martingale,
$P(R_n=k)=\frac1{n+1}$ for $k=1,\dots,n+1$, and $M_n$ converges almost surely to a limit $M_\infty$ that is uniform on $[0,1]$.
\end{theorem}
\begin{proof}
Given the past, a red ball is added with probability $M_n$, so $\mathbb E[R_{n+1}\mid\mathcal F_n]=R_n+M_n=R_n\frac{n+3}{n+2}$, and $\mathbb E[M_{n+1}\mid\mathcal F_n]=M_n$ (rule (e) and the
formula of Section~1, the past being discrete). A given sequence of $n$ draws with $j$ red ones has probability
$\frac{j!\,(n-j)!}{(n+1)!}$: the numerators of the successive factors are $1,2,\dots,j$ for the red draws and $1,\dots,n-j$ for the blue ones, in some order,
and the denominators are $2,3,\dots,n+1$. Multiplying by the $\binom nj$ orders, $P(R_n=j+1)=\frac1{n+1}$.

$M_n$ is a non-negative martingale, so it converges almost surely (Theorem~\ref{thm:supconv}), hence in law (Proposition~\ref{prop:modes}). Now $P(M_n\le x)$ is
the number of $k\in\{1,\dots,n+1\}$ with $k\le x(n+2)$, divided by $n+1$, which tends to $x$ for $x\in[0,1]$. By Proposition~\ref{prop:law}, the distribution
function of $M_\infty$ equals $x$ at its continuity points, which are dense; by right-continuity it equals $x$ on $[0,1]$.
\end{proof}

\begin{figure}[ht]
\centering
\fig{ch10-polya}
\caption{Twenty runs of P\'olya's urn: the proportion of red balls over $1000$ draws. Each run settles down, but at a
random level: the urn remembers its first draws forever.}
\label{fig:polya}
\end{figure}

The limit is random: early luck is reinforced and never forgotten, a model for the success of technologies or of
species that grow by imitation. The same theorem keeps a promise of Chapter~3 about branching processes. Recall
that $Z_0=1$ and $Z_{n+1}=\sum_{i\le Z_n}\xi_{n+1,i}$, where the $\xi_{n,i}$ are independent with the offspring law $(p_k)$, of mean $m$ and generating function $G$.

\begin{theorem}[growth of a branching process]\label{thm:gwlimit}
Let $0<m<\infty$ and $W_n=\frac{Z_n}{m^n}$.
\begin{enumerate}[label=(\alph*)]
\item $W_n$ is a non-negative martingale, which converges almost surely to a limit $W$.
\item If $m\le1$ and $p_1<1$, then $W=0$ almost surely.
\item If $m>1$ and the offspring law has a finite variance $\sigma^2$, then $W_n\to W$ in $L^2$, $\mathbb E[W]=1$, $\mathrm{Var}(W)=\frac{\sigma^2}{m(m-1)}$, and up to a null set
$\{W=0\}$ is the event of extinction.
\end{enumerate}
\end{theorem}
\begin{proof}
(a) Let $\mathcal F_n$ be generated by the $\xi_{k,i}$ with $k\le n$. For $A\in\mathcal F_n$, split according to the value of $Z_n$; as the $\xi_{n+1,i}$ are independent of
$\mathcal F_n$,
\[
\mathbb E[Z_{n+1}\mathbf 1_A]=\sum_kP(A,Z_n=k)\,km=\mathbb E[mZ_n\mathbf 1_A].
\]
So $\mathbb E[Z_{n+1}\mid\mathcal F_n]=mZ_n$, which gives the martingale property, and $\mathbb E[Z_n]=m^n$. Theorem~\ref{thm:supconv} gives convergence.

(b) By Theorem~\ref{thm:extinction} the population dies out almost surely, so $W_n=0$ for large $n$.

(c) In the same way, since $\mathbb E[(\sum_{i\le k}\xi_i)^2]=k\sigma^2+k^2m^2$, we get $\mathbb E[Z_{n+1}^2]=\sigma^2m^n+m^2\,\mathbb E[Z_n^2]$, so
$\mathbb E[W_{n+1}^2]=\mathbb E[W_n^2]+\sigma^2m^{-n-2}$ and
\[
\mathbb E[W_n^2]=1+\sigma^2\sum_{k=0}^{n-1}m^{-k-2}\le1+\frac{\sigma^2}{m(m-1)}.
\]
By Theorem~\ref{thm:L2conv}, $W_n\to W$ in $L^2$, so $\mathbb E[W]=\lim\mathbb E[W_n]=1$ and $\mathbb E[W^2]=\lim\mathbb E[W_n^2]=1+\frac{\sigma^2}{m(m-1)}$.

Let $r=P(W=0)$. The descendants of the $i$-th child of the ancestor form a Galton--Watson process $Z^{(i)}$, with
$Z_{n+1}=\sum_{i\le Z_1}Z^{(i)}_n$; these processes are independent, with the law of $Z$, and independent of $Z_1$, because they are built
from disjoint families of the $\xi$'s. Dividing by $m^{n+1}$ and letting $n\to\infty$, $W=\frac1m\sum_{i\le Z_1}W^{(i)}$, with $W^{(i)}$ independent copies of $W$,
independent of $Z_1$. So $W=0$ if and only if all $W^{(i)}$, $i\le Z_1$, vanish, and $r=\sum_kp_kr^k=G(r)$. Since $m>1$, some $p_k$ with $k\ge2$ is positive, so
$G$ is strictly convex on $[0,1]$ and $G(s)=s$ has at most two solutions there: the extinction probability $q$ and $1$ (Theorem~\ref{thm:extinction}). As
$\mathbb E[W]=1$, $r<1$, so $r=q$. Extinction implies $W=0$, and both events have probability $q$, so they differ by a null set.
\end{proof}

\begin{remark}
So a surviving population grows like $Wm^n$ with $W>0$: the naive mean is right, up to a random factor fixed by the early
generations. Kesten and Stigum (1966) proved that the condition $\sigma^2<\infty$ can be weakened to $\mathbb E[\xi\ln\xi]<\infty$, and that $W=0$ almost surely
otherwise; a short proof by Lyons, Pemantle and Peres is listed in the sources.
\end{remark}

\section{Azuma's inequality}

Hoeffding's inequality (Theorem~\ref{thm:hoeffding}) bounds the deviations of a sum of independent bounded variables. Its proof only used the
fact that each new term has mean $0$ given the past, so it extends to martingales.

\begin{lemma}\label{lem:condhoeffding}
Let $D$ be a random variable with $|D|\le c$ and $\mathbb E[D\mid\mathcal G]=0$. Then $\mathbb E[e^{\lambda D}\mid\mathcal G]\le e^{\lambda^2c^2/2}$ for every real $\lambda$.
\end{lemma}
\begin{proof}
We may assume $c>0$. By convexity, on $[-c,c]$ the function $x\mapsto e^{\lambda x}$ lies below its chord:
\[
\begin{aligned}e^{\lambda x}&\le\frac{c-x}{2c}e^{-\lambda c}+\frac{c+x}{2c}e^{\lambda c}\eqbr=\cosh(\lambda c)+\frac xc\sinh(\lambda c).\end{aligned}
\]
Put $x=D$ and take conditional expectations (rules (a), (b), (d)): $\mathbb E[e^{\lambda D}\mid\mathcal G]\le\cosh(\lambda c)$. Finally $\cosh y=\sum_k\frac{y^{2k}}{(2k)!}\le\sum_k\frac{y^{2k}}{2^kk!}=e^{y^2/2}$,
since $(2k)!=\prod_{j=1}^k(2j-1)(2j)\ge\prod_{j=1}^k2j=2^kk!$.
\end{proof}

\begin{theorem}[Azuma--Hoeffding]\label{thm:azuma}
Let $M$ be a martingale with $|M_k-M_{k-1}|\le c_k$ for constants $c_k$, and $C=\sum_{k\le n}c_k^2$. For $t>0$,
\[
P(M_n-M_0\ge t)\le e^{-t^2/2C},
\]
and the same bound holds for $P(M_n-M_0\le-t)$.
\end{theorem}
\begin{proof}
Let $\lambda>0$ and $Y_n=e^{\lambda(M_n-M_0)}$. As $|M_{n-1}-M_0|\le\sum_kc_k$, $Y_{n-1}$ is bounded and $\mathcal F_{n-1}$-measurable, so by rules (c) and (e) and the lemma,
\[
\begin{aligned}\mathbb E[Y_n]&=\mathbb E\bigl[Y_{n-1}\,\mathbb E[e^{\lambda D_n}\mid\mathcal F_{n-1}]\bigr]\eqbr\le e^{\lambda^2c_n^2/2}\,\mathbb E[Y_{n-1}],\end{aligned}
\]
and by induction $\mathbb E[Y_n]\le e^{\lambda^2C/2}$. Markov's inequality gives $P(M_n-M_0\ge t)\le e^{-\lambda t+\lambda^2C/2}$; take $\lambda=\frac tC$. For the lower tail, apply this to $-M$.
\end{proof}

For a sum of independent variables with $|X_k-\mathbb E X_k|\le c_k$, this is exactly Hoeffding's bound with intervals of length $2c_k$. The gain is that
$M$ may be any martingale. The most useful ones are the Doob martingales of Example~\ref{ex:martingales}(d), which reveal the inputs of a
complicated function one at a time.

\begin{theorem}[bounded differences]\label{thm:bounded}
Let $X_1,\dots,X_n$ be independent, $X_i$ taking values in a countable set $S_i$, and $f:S_1\times\dots\times S_n\to\R$ such that $|f(x)-f(x')|\le c_i$ whenever $x$ and $x'$
differ only in the $i$-th coordinate. Let $F=f(X_1,\dots,X_n)$ and $C=\sum_ic_i^2$. For $t>0$,
\[
P\bigl(|F-\mathbb E[F]|\ge t\bigr)\le2e^{-t^2/2C}.
\]
\end{theorem}
\begin{proof}
Changing the coordinates one at a time, $|f(x)-f(y)|\le\sum_ic_i$ for all $x,y$, so $f$ is bounded. Let $\mathcal F_k=\sigma(X_1,\dots,X_k)$ and $M_k=\mathbb E[F\mid\mathcal F_k]$, so that
$M_0=\mathbb E[F]$ and $M_n=F$. The $\sigma$-algebra $\mathcal F_k$ comes from the countable partition into the events $\{(X_1,\dots,X_k)=x\}$, so by Section~1 and
independence, $M_k=g_k(X_1,\dots,X_k)$ where
\[
g_k(x_1,\dots,x_k)=\mathbb E\bigl[f(x_1,\dots,x_k,X_{k+1},\dots,X_n)\bigr].
\]
Write $x'=(x_1,\dots,x_{k-1})$. By independence again, $g_{k-1}(x')=\sum_zP(X_k=z)\,g_k(x',z)$, and $|g_k(x',x_k)-g_k(x',z)|\le c_k$ since the
same holds for $f$ inside the expectation. So $|M_k-M_{k-1}|\le\sum_zP(X_k=z)c_k=c_k$, and Theorem~\ref{thm:azuma}, applied to both tails, concludes.
\end{proof}

For example, throw $m$ balls independently into $n$ bins and let $F$ be the number of empty bins. Moving one ball changes $F$ by at most
$1$, so $P(|F-\mathbb E F|\ge t)\le2e^{-t^2/2m}$, although the bins are far from independent. With $m=n=10^6$, $\mathbb E[F]=n(1-\frac1n)^m\approx367\,879$, and the
probability of a deviation of $5000$ is at most $2e^{-12.5}<10^{-5}$. The true standard deviation is about $312$: Azuma loses a constant
factor, but finds the right scale $\sqrt m$ without knowing anything about the dependence. McDiarmid (1989) improved the exponent to
$\frac{2t^2}C$, which matches Hoeffding exactly.

\section{Exercises}

\begin{exercise}[$\star$]\label{ex:ruinasym}
\emph{Gambler's ruin, unfair case.} Let $S$ be the walk with steps $+1$ with probability $p\ne\frac12$ and $-1$ with probability $q$, started at $a\in\{0,\dots,N\}$, and $T$ its
exit time from $\{1,\dots,N-1\}$. Using the martingales $(\frac qp)^{S_n}$ and $S_n-(p-q)n$, find $P(S_T=N)$ and $\mathbb E[T]$, and compare with Theorem~\ref{thm:ruin}.
\end{exercise}

\begin{exercise}[$\star$]\label{ex:patterns}
A fair coin is tossed repeatedly. Compute the expected waiting times for HH, HT, HTH and HHT. Why does HH take longer than HT,
although both patterns occur at a quarter of all positions in the long run?
\end{exercise}

\begin{exercise}[$\star$]\label{ex:doubling}
\emph{The doubling strategy.} A player bets $1$ euro on a fair coin, doubles the stake after each loss, and stops at the first win.
Show that the net gain at the stopping time $T$ is $1$ almost surely. Compute the expected loss accumulated just before the
winning round, and say which hypothesis of Theorem~\ref{thm:ost} fails for the fortune.
\end{exercise}

\begin{exercise}[$\star\star$]\label{ex:wald2}
\emph{Wald's second identity.} Let $\xi_1,\xi_2,\ldots$ be independent with the same law, of mean $0$ and variance $\sigma^2$, $S_n=\xi_1+\dots+\xi_n$, and $T$ a stopping time
with $\mathbb E[T]<\infty$. Show that $\mathbb E[S_T]=0$ and $\mathbb E[S_T^2]=\sigma^2\,\mathbb E[T]$. (Apply Theorem~\ref{thm:L2conv} to the stopped walk.)
\end{exercise}

\begin{exercise}[$\star\star$]\label{ex:polyabeta}
\emph{P\'olya's urn in general.} The urn starts with $a$ red and $b$ blue balls ($a,b\ge1$ integers). Show that the proportion of red balls is a
martingale, that the number $K$ of red balls drawn in the first $n$ steps satisfies
\[
P(K=k)=\frac{\binom{a+k-1}k\binom{b+n-k-1}{n-k}}{\binom{a+b+n-1}n},
\]
and that the limit proportion has the Beta density $\frac{x^{a-1}(1-x)^{b-1}}{B(a,b)}$ on $[0,1]$, where $B(a,b)=\frac{(a-1)!\,(b-1)!}{(a+b-1)!}$.
\end{exercise}

\begin{exercise}[$\star$]\label{ex:signs}
Let $\varepsilon_n$ be independent fair signs. Show that $\sum_n\varepsilon_nn^{-\alpha}$ converges almost surely for $\alpha>\frac12$, and that for $\alpha\le\frac12$ the partial sums have a
variance tending to infinity. Compute the variance of $\sum_n\frac{\varepsilon_n}n$.
\end{exercise}

\begin{exercise}[$\star\star$]\label{ex:chromatic}
\emph{The chromatic number of a random graph.} In the random graph $G(n,\frac12)$, each of the $\binom n2$ possible edges between $n$ vertices is
present independently with probability $\frac12$. Its chromatic number $\chi$ is the least number of colours that colour the vertices so
that adjacent vertices differ. Show that $P(|\chi-\mathbb E\chi|\ge\lambda\sqrt n)\le2e^{-\lambda^2/2}$ for every $\lambda>0$, although $\mathbb E\chi$ itself is hard to compute. (Reveal the vertices one at a time.)
\end{exercise}

\begin{exercise}[$\star\star$]\label{ex:likelihood}
\emph{Likelihood ratios.} Let $X_1,X_2,\ldots$ be independent with law $p$ on a finite set, and $q$ another law with $q(x)>0$ exactly when $p(x)>0$. Show
that $L_n=\prod_{k\le n}\frac{q(X_k)}{p(X_k)}$ is a martingale, that it converges almost surely, and that the limit is $0$ if $q\ne p$. (Use the strong law for
$\ln\frac{q(X_k)}{p(X_k)}$ and Gibbs' inequality.) A statistician who doubts that the data come from $p$ thus becomes, almost surely, more and
more certain, although the expected likelihood ratio stays equal to $1$.
\end{exercise}

\section{Hints}

\begin{hint}{ex:ruinasym}
Both martingales are bounded on $\{T\ge n\}$, or have bounded increments, and $\mathbb E[T]<\infty$ by Lemma~\ref{lem:exit}. Optional stopping gives two linear
equations.
\end{hint}
\begin{hint}{ex:patterns}
Theorem~\ref{thm:abra} with $m=2$. List the overlaps.
\end{hint}
\begin{hint}{ex:doubling}
If the first win is at round $T$, the losses before it are $1+2+\dots+2^{T-2}$. $T$ is geometric.
\end{hint}
\begin{hint}{ex:wald2}
$S_{T\wedge n}$ is a martingale with orthogonal increments $\xi_k\mathbf 1_{\{T\ge k\}}$, and $\mathbb E[\xi_k^2\mathbf 1_{\{T\ge k\}}]=\sigma^2P(T\ge k)$.
\end{hint}
\begin{hint}{ex:polyabeta}
A sequence of draws with $k$ red ones has a probability that does not depend on the order. For the limit, compare $n\,P(K=k)$ with
the Beta density at $\frac kn$.
\end{hint}
\begin{hint}{ex:signs}
The example after Theorem~\ref{thm:L2conv}. For the variance, $\sum\frac1{n^2}=\frac{\pi^2}6$ (Tome~0).
\end{hint}
\begin{hint}{ex:chromatic}
Let $X_i$ be the set of edges between vertex $i$ and the vertices $1,\dots,i-1$. Changing $X_i$ changes $\chi$ by at most $1$.
\end{hint}
\begin{hint}{ex:likelihood}
The mean of $\frac{q(X)}{p(X)}$ is $\sum_{x:\,p(x)>0}q(x)=1$. The mean of $\ln\frac{q(X)}{p(X)}$ is negative unless $q=p$.
\end{hint}

\section{Solutions}

\begin{solution}{ex:ruinasym}
\idea Each martingale, stopped at the exit time, turns into one equation, and the two unknowns are the probability of
winning and the mean duration.
\writeup Let $r=\frac qp$ and $h=P(S_T=N)$. By Lemma~\ref{lem:exit}, $\mathbb E[T]<\infty$. The martingale $r^{S_n}$ stopped at $T$ is bounded by $\max(1,r^N)$, so by
Theorem~\ref{thm:ost}(b), $r^a=\mathbb E[r^{S_T}]=hr^N+(1-h)$, and $h=\frac{r^a-1}{r^N-1}$. The martingale $S_n-(p-q)n$ has increments bounded by $2$, so by Theorem~\ref{thm:ost}(c),
$a=\mathbb E[S_T]-(p-q)\mathbb E[T]=Nh-(p-q)\mathbb E[T]$, and $\mathbb E[T]=\frac{a-Nh}{q-p}$. These are the formulas of Theorem~\ref{thm:ruin}, obtained without a recurrence.
\end{solution}

\begin{solution}{ex:patterns}
\idea A pattern that overlaps itself wastes part of its occurrences in clumps, so the gaps between clumps are longer.
\writeup By Theorem~\ref{thm:abra} with $m=2$: HH has overlaps $1,2$, so $2+4=6$ tosses; HT has only the overlap $2$: $4$ tosses; HTH has overlaps $1,3$:
$2+8=10$; HHT has only $3$: $8$. Both HH and HT occur at a quarter of all positions, so in a long sequence of length $L$ there are about
$\frac L4$ occurrences of each. But occurrences of HH come in runs (HHHH contains three), so they gather into fewer clumps,
separated by longer gaps; the waiting time from the start resembles a gap between clumps. Occurrences of HT never
overlap, and are evenly spread.
\end{solution}

\begin{solution}{ex:doubling}
\idea The strategy is a martingale transform with unbounded stakes, played for an unbounded time: optional stopping fails
because the fortune is unbounded below.
\writeup If the first win is at round $T$, the player has lost $1+2+\dots+2^{T-2}=2^{T-1}-1$ and then wins $2^{T-1}$: net gain $1$. $T$ is geometric with
$P(T=k)=2^{-k}$, finite almost surely. The loss just before the winning round is $2^{T-1}-1$, with expectation $\sum_k2^{-k}(2^{k-1}-1)=\infty$. The fortune
$X_n$ is a martingale (Proposition~\ref{prop:transform} applies at each finite time, since the stake $2^{k-1}$ at round $k$ is bounded for fixed $k$), and $\mathbb E[X_T]=1\ne0$.
Here $\mathbb E[T]=2<\infty$, but the increments are not bounded, so (c) fails; and $X_{T\wedge n}$ is not bounded below, so (b) fails. A player with a finite
credit line $2^L-1$ is ruined with probability $2^{-L}$, losing $2^L-1$, which exactly balances the gain $1$ in the other cases.
\end{solution}

\begin{solution}{ex:wald2}
\idea The stopped walk is a martingale whose increments are orthogonal, and their squares add up to $\sigma^2\,\mathbb E[T]$; boundedness in $L^2$
then lets us pass to the limit.
\writeup $S_{T\wedge n}=\sum_{k\le n}\xi_k\mathbf 1_{\{T\ge k\}}$ is a martingale (Proposition~\ref{prop:stopped}), in $L^2$. The event $\{T\ge k\}$ is in $\mathcal F_{k-1}$, hence independent of $\xi_k$,
so $\mathbb E[\xi_k^2\mathbf 1_{\{T\ge k\}}]=\sigma^2P(T\ge k)$. By Lemma~\ref{lem:orth},
\[
\begin{aligned}\mathbb E[S_{T\wedge n}^2]&=\sigma^2\sum_{k\le n}P(T\ge k)\eqbr=\sigma^2\,\mathbb E[T\wedge n]\le\sigma^2\,\mathbb E[T].\end{aligned}
\]
So the stopped walk is bounded in $L^2$; by Theorem~\ref{thm:L2conv} it converges almost surely and in $L^2$, and since $T<\infty$ almost surely its limit is
$S_T$. Convergence in $L^2$ gives $\mathbb E[S_T]=\lim\mathbb E[S_{T\wedge n}]=0$ and $\mathbb E[S_T^2]=\lim\mathbb E[S_{T\wedge n}^2]=\sigma^2\,\mathbb E[T]$ (monotone convergence on the right). For the simple walk
leaving $]0,N[$ from $a$, this gives again $\mathbb E[T]=\mathbb E[(S_T-a)^2]=a(N-a)$.
\end{solution}

\begin{solution}{ex:polyabeta}
\idea Exchangeability: the probability of a sequence of draws depends only on how many reds it contains, so everything
reduces to rising factorials, and their asymptotics are powers.
\writeup After $n$ steps there are $a+b+n$ balls, of which $R_n$ red; as in Theorem~\ref{thm:polyaurn}, $\mathbb E[R_{n+1}\mid\mathcal F_n]=R_n\frac{a+b+n+1}{a+b+n}$, so $\frac{R_n}{a+b+n}$ is a
martingale. Write $x^{(j)}=x(x+1)\cdots(x+j-1)$ for the rising factorial. A sequence of $n$ draws with $k$ reds has probability
\[
\frac{a^{(k)}\,b^{(n-k)}}{(a+b)^{(n)}},
\]
whatever the order: the red numerators run through $a,\dots,a+k-1$, the blue ones through $b,\dots,b+n-k-1$. Multiplying by $\binom nk$ and writing
the rising factorials with factorials gives the formula. Now $\binom{a+k-1}k=\frac{(k+1)\cdots(k+a-1)}{(a-1)!}$ lies between $\frac{k^{a-1}}{(a-1)!}$ and $\frac{(k+a)^{a-1}}{(a-1)!}$, and similarly for
the other two binomials. So $P(K=k)$ lies between $\frac1{B(a,b)}\frac{k^{a-1}(n-k)^{b-1}}{(n+a+b)^{a+b-1}}$ and $\frac1{B(a,b)}\frac{(k+a)^{a-1}(n-k+b)^{b-1}}{n^{a+b-1}}$. Summing over $k\le xn$, both
bounds are Riemann sums, up to factors tending to $1$ and shifts of order $\frac1n$, of $\int_0^x\frac{u^{a-1}(1-u)^{b-1}}{B(a,b)}\dd u$, which is continuous in $x$. So $P(\frac Kn\le x)$
tends to this integral. The limit proportion $M_\infty$ of red balls is also the limit of $\frac Kn$, since $R_n=a+K$, so it has the Beta law
(Propositions~\ref{prop:modes} and~\ref{prop:law}, as in Theorem~\ref{thm:polyaurn}). The normalisation $\int_0^1u^{a-1}(1-u)^{b-1}\dd u=B(a,b)$ follows: the total mass is $1$.
\end{solution}

\begin{solution}{ex:signs}
\idea Random signs make a martingale whose $L^2$ norm is the sum of the squares of the terms.
\writeup $M_n=\sum_{k\le n}\varepsilon_kk^{-\alpha}$ is a martingale (Example~\ref{ex:martingales}(a)) with $\mathbb E[M_n^2]=\sum_{k\le n}k^{-2\alpha}$, which is bounded if and only if $2\alpha>1$. For $\alpha>\frac12$
Theorem~\ref{thm:L2conv} gives almost sure convergence; for $\alpha\le\frac12$ the variance tends to infinity. For $\alpha=1$ the limit has variance $\sum\frac1{n^2}=\frac{\pi^2}6$. (For
$\alpha\le\frac12$ the series diverges almost surely, since its terms do not tend to $0$ when $\alpha\le0$, and for $0<\alpha\le\frac12$ by Kolmogorov's three-series
theorem, Tome~10.)
\end{solution}

\begin{solution}{ex:chromatic}
\idea Reveal the graph vertex by vertex: each step changes the chromatic number by at most one, so Azuma applies with
$n-1$ steps of size $1$.
\writeup For $2\le i\le n$, let $X_i$ be the set of edges present between vertex $i$ and the vertices $1,\dots,i-1$: the $X_i$ are independent (they involve
disjoint sets of edges), with values in finite sets, and the graph is a function of $(X_2,\dots,X_n)$. If two graphs differ only in the
edges at vertex $i$, their chromatic numbers differ by at most $1$: from a colouring of one, give vertex $i$ a new colour to get a
colouring of the other. Theorem~\ref{thm:bounded} with $n-1$ coordinates and $c_i=1$ gives $P(|\chi-\mathbb E\chi|\ge t)\le2e^{-t^2/2(n-1)}\le2e^{-t^2/2n}$; take $t=\lambda\sqrt n$.
Since $\mathbb E\chi$ is asymptotic to $\frac n{2\log_2n}$ (Bollob\'as, 1988, with martingales again), this says that $\chi$ is concentrated in a window much smaller than its size.
\end{solution}

\begin{solution}{ex:likelihood}
\idea The likelihood ratio is a product of independent factors of mean at most $1$; its logarithm is a random walk with
negative drift.
\writeup Let $Y_k=\frac{q(X_k)}{p(X_k)}$, defined and positive since $p(X_k)>0$ almost surely. The $Y_k$ are independent, positive, with
$\mathbb E[Y_k]=\sum_{x:\,p(x)>0}p(x)\frac{q(x)}{p(x)}=\sum_xq(x)=1$, so $L_n$ is a martingale (Example~\ref{ex:martingales}(c)), non-negative, and it converges almost
surely by Theorem~\ref{thm:supconv}. Now $\ln L_n=\sum_{k\le n}\ln Y_k$ is a sum of independent bounded variables (the set is finite) of mean
$\sum_xp(x)\ln\frac{q(x)}{p(x)}=-D(p\|q)$, which is $<0$ when $q\ne p$ (Gibbs' inequality, Proposition~\ref{prop:gibbs}). By the strong law (Corollary~\ref{cor:slln}),
$\frac1n\ln L_n\to-D(p\|q)$ almost surely, so $L_n\to0$, exponentially fast. Yet $\mathbb E[L_n]=1$ for all $n$: the expectation is carried by rarer and
rarer sequences on which $L_n$ is huge.
\end{solution}

\section*{Sources and further reading}

\begin{itemize}
\item D.~Williams, \emph{Probability with Martingales}, Cambridge University Press, 1991. Conditional expectation as a
projection, upcrossings, the ABRACADABRA problem (an exercise of Chapter~10) and much more, in a lively style.
\item R.~Durrett, \emph{Probability: Theory and Examples}, 5th ed., Cambridge University Press, 2019, Chapter~4.
\item G.~Grimmett and D.~Stirzaker, \emph{Probability and Random Processes}, 4th ed., Oxford University Press, 2020,
Chapter~12.
\item J.~L.~Doob, \emph{Stochastic Processes}, Wiley, 1953. The book that made martingales a theory.
\item S.-Y.~R.~Li, \emph{A martingale approach to the study of occurrence of sequence patterns in repeated experiments},
Annals of Probability 8 (1980), 1171--1176.
\item K.~Azuma, \emph{Weighted sums of certain dependent random variables}, T\^ohoku Mathematical Journal 19 (1967),
357--367.
\item C.~McDiarmid, \emph{On the method of bounded differences}, in Surveys in Combinatorics 1989, London Mathematical
Society Lecture Note Series 141, Cambridge University Press, 1989, 148--188.
\item E.~Shamir and J.~Spencer, \emph{Sharp concentration of the chromatic number on random graphs $G_{n,p}$},
Combinatorica 7 (1987), 121--129.
\item B.~Bollob\'as, \emph{The chromatic number of random graphs}, Combinatorica 8 (1988), 49--55.
\item R.~Lyons, R.~Pemantle and Y.~Peres, \emph{Conceptual proofs of $L\log L$ criteria for mean behavior of branching
processes}, Annals of Probability 23 (1995), 1125--1138.
\item M.~Mitzenmacher and E.~Upfal, \emph{Probability and Computing}, 2nd ed., Cambridge University Press, 2017,
Chapter~13: martingales and Azuma's inequality for algorithms.
\end{itemize}
